\section{Data}
\label{sec:data}

Our data curation process consisted of acquisition, processing, filtering, and captioning, as illustrated in Fig. \ref{fig:data_pipeline}. We first describe our data acquisition, processing, and filtering in \cref{subsec:data_source}. Then we present our motion captioning process in \cref{subsec:captioning} and motion taxonomy in \cref{subsec:stats}.

\begin{figure}[h!]
    \centering
    \includegraphics[width=\linewidth, trim=30 150 100 150, clip]{figures/data_engine_pipeline.pdf}
    \vspace{-10pt}
    \caption{Overview of the data processing pipeline.}
    \label{fig:data_pipeline}
\end{figure}




\subsection{Data Acquisition and Filtering}
\label{subsec:data_source}


\paragraph{Data Acquisition.}



Our dataset was constructed from three complementary data sources: in-the-wild human motion videos, motion capture data, and 3D animation assets.
Human motion videos provided rich and diverse actions captured across varied real-world scenarios.
Motion capture data, typically recorded in controlled indoor environments, offered high-quality motion data but were limited in scene diversity.
3D animation assets, hand-crafted by professional artists for game production, exhibited exceptional motion quality but were relatively expensive and limited in quantity.


For video data, we started with a massive collection of 12 million high-quality, in-the-wild video clips from HunyuanVideo~\citep{hunyuanvideo2025}.
We first processed these clips through a rigorous pre-processing stage, employing shot boundary detection to segment videos into coherent scenes and a human detector to search for clips containing human subjects.
Subsequently, for the resulting candidate clips, we utilized a state-of-the-art human motion extraction algorithm GVHMR~\citep{shen2024gvhmr}, to reconstruct 3D human tracks of SMPL-X~\citep{smplx} parameters.
For motion capture data and 3D animation assets, we acquired approximately 500 hours of motion sequences.












\paragraph{Motion Processing and Filtering}

To facilitate downstream model training, we standardized all motion data onto a unified SMPL-H skeleton. We processed all motions in three steps: retargeting, low-quality motion filtering, and canonicalization.
We retargeted all motions onto a neutral SMPL-H skeleton~\citep{smplh}.
For motion data originally in SMPL~\citep{smpl}, SMPL-H~\citep{smplh}, or SMPL-X~\citep{smplx} formats, we employed mesh fitting to convert them into the unified SMPL-H representation.
For data with other skeletal structures, we applied a retargeting tool to map them to the SMPL-H skeleton.
For all retargeted data, we established a comprehensive filtering pipeline to remove low-quality motion clips.
This included removal of duplicates, abnormal poses, outliers of joint velocity, detection of anomalous displacements, pruning of static motions, and detection of artifacts such as foot-sliding.
After filtering, we applied canonicalization to standardize the data.
All motions were resampled to 30 fps, and sequences longer than 12 seconds were segmented into multiple clips.
Each motion was normalized to a canonical coordinate frame: Y-axis up, the starting position centered at the origin, the lowest body point aligned on the ground plane, and initial facing direction along the positive Z-axis.
Ultimately, we obtained over 3000 hours of motion data, including 400 hours of high-quality 3D motion data.




\subsection{Motion Captioning}
\label{subsec:captioning}





Our captioning pipeline primarily leveraged vision-language models (VLMs) for automatic annotation, followed by manual verification and refinement. Finally, we employed large language models (LLMs) to diversify and enrich the textual descriptions.

For video-sourced data, we directly used the original videos corresponding to each motion sequence.
For 3D motion data, we textured and rendered the SMPL-H model to generate synthetic videos.
These videos, along with optimized prompts dedicated to human motion, were fed into a VLM (\eg, Gemini-2.5-Pro) to obtain preliminary captions and action keywords.
To obtain as accurate a text-motion pair as possible for the high-quality data (i.e., rendered motions),
we conducted manual verification on the VLM outputs, correcting erroneous descriptions and supplementing missing key motion components.

For all generated captions, we utilized an LLM to standardize the caption structure while preserving the original semantics and to create diverse paraphrases for data augmentation.












\subsection{Taxonomy}
\label{subsec:stats}


To organize our dataset, we established a three-level hierarchical motion taxonomy based on textual captions and keywords. As illustrated in Figure \ref{fig:data_overview}, the hierarchy began with six coarse-grained, top-level categories: (a) Locomotion, (b) Sports \& Athletics, (c) Fitness \& Outdoor Activities, (d) Daily Activities, (e) Social Interactions \& Leisure, and (f) Game Character Actions. These categories were progressively subdivided, culminating in over 200 fine-grained motion classes at the leaf level.


\begin{figure*}[t]
    \centering
    \begin{minipage}[b]{\textwidth}
        \centering
        \includegraphics[width=\textwidth, trim=0 20 0 20]{figures/data/category_sunburst.pdf}
    \end{minipage}
    \hfill
    \caption{The hierarchy of our motion categories.}
    \label{fig:data_overview}
\end{figure*}











